\begin{lemma}[Proved: the dual-number Bockstein detects every nonzero primary factor]
\label{C:epsilon_Bockstein}
Let $A$ be a discrete valuation ring of characteristic zero whose
maximal ideal is $(2)$, put $F=A/2A$, and let
$k=A[\epsilon]/\epsilon^2$. For a finite free $k$-module $M$ with an
involution, write $L=M/\epsilon M$ and
\[
 V(M)=\widehat H^0(C_2,L),\qquad
 W(M)=\widehat H^1(C_2,L).
\]
These are finite-dimensional $F$-spaces. The connecting maps of
$0\to L\xrightarrow{\epsilon}M\to L\to0$, regarded as a sequence of
$A[C_2]$-modules, give natural maps
\[
 b^0:V(M)\longrightarrow W(M),\qquad
 b^1:W(M)\longrightarrow V(M).
\]
Let $u=b^1b^0$ on $V(M)$ and $u=b^0b^1$ on $W(M)$.
For every monic irreducible $q\in F[u]$ other than $u$, the function
\[
 \chi_q^\epsilon(M)
   =\operatorname{length}_{F[u]}V(M)_q\pmod2               \tag{C.1.9}
\]
is additive on admissible short exact sequences and hence defines
$R_k(C_2)\to\mathbb Z/2$.
\end{lemma}
\begin{proof}
Use the two-periodic Tate complex whose alternating differentials
are $T-1$ and $T+1$. Its cohomology is annihilated by $2$, since their
sum is $2$, and is finite over $A$. Thus it is finite-dimensional
over $F$. The sequence defining $b^i$ is exact because $M$ is free
over $k$, and its terms are finite free over $A$.

For clarity, the connecting maps do not necessarily satisfy
$b^1b^0=0$. Choose an $A$-linear section $L\to M$, and write
\[
 T=T_0+\epsilon U,\qquad T_0^2=1,\quad T_0U+UT_0=0.
\]
If $x$ is a Tate cycle in $L$, its lift has differential
$\epsilon Ux$; thus $b^i[x]=[Ux]$. This also follows directly
from the connecting-map definition and makes its independence of
the chosen section explicit.

An admissible sequence $0\to M_A\to M_E\to M_B\to0$ remains
exact after reduction modulo $\epsilon$, because it splits as a
sequence of $k$-modules. Its Tate sequence is the exact six-cycle
\[
 V_A\longrightarrow V_E\longrightarrow V_B
 \xrightarrow{\partial^0}W_A\longrightarrow W_E
 \longrightarrow W_B\xrightarrow{\partial^1}V_A.       \tag{C.1.10}
\]
The Bocksteins commute with the first two maps and anticommute
with the connecting maps. Here is a direct check of the latter
assertion, including its hypotheses. Choose an underlying $k$-linear
splitting and write the two-periodic differential of the middle
complex in block form
\[
 d_E^i=\begin{pmatrix}d_A^i&X\\0&d_B^i\end{pmatrix},
 \quad d_A^i=d_{A,0}^i+\epsilon U_A,\quad
 d_B^i=d_{B,0}^i+\epsilon U_B,\quad X=X_0+\epsilon X_1.
\]
The coefficient of $\epsilon$ in the off-diagonal block of
$d_E^{i+1}d_E^i=0$ is
\[
 U_AX_0+d_{A,0}^{i+1}X_1+X_1d_{B,0}^i+X_0U_B=0.
\]
On a $d_{B,0}^i$-cycle the middle two terms give a boundary.
Since $\partial^i[x]=[X_0x]$, this proves
$b_A^{i+1}\partial^i=-\partial^{i+1}b_B^i$.
In particular $u$ commutes with every map in (C.1.10).

Take the $q$-primary parts. Primary decomposition is exact for
finite-dimensional $F[u]$-modules, so the resulting six-cycle is
exact. Since $q\ne u$, the operator $u$ is invertible on these
parts, and both $b^0$ and $b^1$ are isomorphisms between the
$V$- and $W$-parts. Using $(b_A^0)^{-1}$ to identify $W_{A,q}$
with $V_{A,q}$ turns the first half of (C.1.10) into an exact
three-cycle
\[
 V_{A,q}\longrightarrow V_{E,q}\longrightarrow V_{B,q}
 \xrightarrow{(b_A^0)^{-1}\partial^0}V_{A,q}.
\]
Exactness at its last vertex follows from the commutation of $b^0$
with $V_A\to V_E$ and $W_A\to W_E$. Every map here is
$F[u]$-linear. If its three image lengths are $r,s,t$, the space
lengths are $t+r,r+s,s+t$, respectively. Their alternating sum
is even. This proves (C.1.9).
\end{proof}

\begin{theorem}[Proved: the entire degree-two group over flat dual numbers]
\label{C:dualnumber_C2_ring}
Keep the assumptions of Lemma~\ref{C:epsilon_Bockstein}; in
particular $A=\mathbb Z_{(2)}$ and $A=\mathbb Z_2$ are allowed.
Let $\mathcal I_F$ be the set of monic irreducible polynomials in
$F[u]$ other than $u$. For $q\in\mathcal I_F$ of degree $d$, lift
the entries of its companion matrix to an invertible matrix
$B_q\in M_d(A)$ and set
\[
 T_q=\begin{pmatrix}I_d&\epsilon B_q\\
                     \epsilon I_d&-I_d\end{pmatrix},
 \quad M_q=(k^{2d},T_q),\qquad
 \eta_q=[M_q]-d(1+s),\quad s=[\mathrm{sgn}].             \tag{C.1.11}
\]
Then, as an abelian group,
\[
 R_k(C_2)=\mathbb Z\{1,s\}\ \oplus
             \bigoplus_{q\in\mathcal I_F}(\mathbb Z/2)\eta_q.
                                                               \tag{C.1.12}
\]
For the internal tensor product on this representation group,
the complete multiplication rules are
\[
 s^2=1,\qquad s\eta_q=\eta_q,\qquad
 \eta_q\eta_{q'}=0\quad(q,q'\in\mathcal I_F).           \tag{C.1.13}
\]
The second summand is precisely the kernel of generic scalar
extension to $R_{k[1/2]}(C_2)=\mathbb Z\{1,s\}$ and also
the cokernel of the degree-two polynomial induction classes.
Thus this is a computation of the full degree-two obstruction,
not merely an injection into that obstruction. The assertion
concerns this family of base rings; it does not settle C.1 for
every commutative ring.
\end{theorem}
\begin{proof}
\emph{The integral starting point.}
We first give an elementary decomposition of every finite free
$A$-module $L$ with an involution $T_0$. Put
$L_\pm=\ker(T_0\mp1)$. These are finite free over the DVR $A$,
their direct sum has full rank in $L$, and $2L\subset L_++L_-$.
The quotient $Q=L/(L_++L_-)$ is an $F$-space, and the maps
\[
 Q\longrightarrow L_+/2L_+,\ [v]\longmapsto[v+T_0v],
 \qquad
 Q\longrightarrow L_-/2L_-,\ [v]\longmapsto[v-T_0v]
\]
are injective. For example if $v+T_0v=2a$ with $a\in L_+$,
then $v-a\in L_-$; the other assertion is identical.
Choose an $F$-basis $[v_1],\ldots,[v_r]$ of $Q$.
The elements $e_i=v_i+T_0v_i$ extend to an $A$-basis of $L_+$,
and the elements $f_i=v_i-T_0v_i$ extend to one of $L_-$,
because their reductions are independent. The relations
$2v_i=e_i+f_i$ show that $L$ has an $A$-basis consisting of
$v_i,T_0v_i$ for $1\le i\le r$, followed by the unused basis
vectors of $L_+$ and $L_-$. On each pair $T_0$ swaps the two
vectors. We have proved
\[
 L\simeq A[C_2]^r\oplus A^a\oplus A_{\mathrm{sgn}}^b.  \tag{C.1.14}
\]

\emph{Lifting and removing the regular summands.}
For an arbitrary object $M$ over $k$, apply (C.1.14) to
$L=M/\epsilon M$. Lift each regular generator $v_i$ to $M$.
The vectors $v_i,Tv_i$, together with lifts of the remaining
basis of $L$, form a $k$-basis of $M$: their change-of-basis
matrix is invertible modulo the nilpotent ideal $(\epsilon)$.
They therefore define a $k[C_2]^r$-submodule with a free
underlying quotient. It splits equivariantly. Indeed, choose
$k$-linear functionals $a_i$ taking values $1,0$ on
$v_i,Tv_i$ and zero on the other basis vectors. The maps
\[
 m\longmapsto a_i(m)+g\,a_i(Tm)\in k[C_2]
\]
are equivariant and jointly retract this submodule. Modulo
$\epsilon$ this is the original regular projection, so the
complement has reduction $A^a\oplus A_{\mathrm{sgn}}^b$.
Regular summands contribute only $1+s$ to $R_k(C_2)$.

\emph{A finite-length matrix problem gives all generators.}
On the complement, lift a basis of its plus and minus reduction.
Its involution has the form
\[
 T=\begin{pmatrix}I_a&0\\0&-I_b\end{pmatrix}
    +\epsilon\begin{pmatrix}0&B\\ C&0\end{pmatrix}.    \tag{C.1.15}
\]
Indeed $T^2=1$ forces the two diagonal blocks of the first-order
term to be zero, since $A$ has no $2$-torsion.
Conjugation by $I+\epsilon X$ changes $B$ and $C$ independently
by arbitrary matrices divisible by $2$. Consequently only
their reductions
\[
 \overline C:F^a\longrightarrow F^b,\qquad
 \overline B:F^b\longrightarrow F^a
\]
matter for the following filtration argument. An invertible
change of basis at the two $F$-spaces lifts to one over $A$.
For a subrepresentation of this two-arrow cycle choose bases
adapted to its two subspaces. The matrix blocks which must vanish
for invariance are divisible by $2$; an $I+\epsilon X$
conjugation makes them exactly zero. The lifted coordinate
subspaces therefore give an invariant free $k$-submodule with
free quotient. Repeating along a composition series expresses
$[M]$ as a sum of the lifts of the simple representations of
this finite-dimensional quiver representation.

Its simple representations can be classified directly. There
are the two one-dimensional vertex representations. If both
vertex spaces of a simple representation are nonzero, the
kernel of either arrow would give a subrepresentation supported
at one vertex; hence both arrows are injective and thus are
isomorphisms. Identify one arrow with the identity. Simplicity
then says that the other is a simple invertible $F[u]$-module,
so it is the companion matrix of one uniquely determined
$q\in\mathcal I_F$. Lifting its basis change and adjusting even
entries in (C.1.15) gives precisely $M_q$. We have proved that
$1,s$ and the $[M_q]$, hence $1,s,\eta_q$, generate the group.
This uses finite-length filtrations to compute the group; it
does not claim that all indecomposable dual-number lattices
have been classified.

\emph{All relations and the generic kernel.}
The determinant of $B_q$ is a unit since $q(0)\ne0$.
The invertible matrix
\[
 J_q=\begin{pmatrix}0&B_q\\-I_d&0\end{pmatrix}
\]
satisfies $J_qT_q=(-T_q)J_q$. Thus $s[M_q]=[M_q]$.
The evaluation sequence (C.1.7), of rank $2d$ here, gives
$2[M_q]=2d(1+s)$, and therefore $2\eta_q=0$ and
$s\eta_q=\eta_q$.

The reduction of $M_q$ modulo $\epsilon$ is the direct sum of
$d$ plus and $d$ minus lines. Its Tate groups are respectively
$F^d$ and $F^d$, and its Bocksteins are
\[
 b^0=I_d,\qquad b^1=\overline B_q.
\]
It follows that
$\chi_{q'}^\epsilon(\eta_q)=1$ if $q'=q$ and is zero otherwise.
The two constant lines have zero Bockstein, so every such
invariant vanishes on them. Generically $(1\pm T_q)/2$ split
$M_q$ into free plus and minus modules of rank $d$: the ring
$k[1/2]=\operatorname{Frac}(A)[\epsilon]/\epsilon^2$ is local.
Thus the generic comparison records the two independent integer
coefficients of $1,s$, and the homomorphisms
$\chi_q^\epsilon$ record independently every proposed
$\mathbb Z/2$ coefficient. These detectors, together with the
generation already proved, give exactly (C.1.12) and its asserted
generic kernel. The degree-two induction image is the span of
$1$ and the regular class $1+s$, which is the same integer
summand.

\emph{Multiplication.}
Let $B=B_q$ and $C=B_{q'}$ have sizes $d,e$. Reduce
$M_q\otimes_k M_{q'}$ modulo $\epsilon$. Its plus Tate space
is the sum of the plus--plus and minus--minus tensor spaces,
each of dimension $de$. The first-order involution is
$U_B\otimes T_{C,0}+T_{B,0}\otimes U_C$. Modulo $2$ its square
on the plus Tate space is the direct sum of two copies of
\[
 \overline B\otimes I_e+I_d\otimes\overline C.          \tag{C.1.16}
\]
This follows by squaring: the two cross terms cancel in
characteristic two and $U_B^2$ acts by $B$ on each sign space,
and similarly for $C$. Hence every primary length is even and
every $\chi_{q''}^\epsilon$ of this tensor product is zero.
Its generic class is $de(1+s)^2=2de(1+s)$; by (C.1.12),
\[
 [M_q][M_{q'}]=2de(1+s).
\]
Since $(1+s)\eta_q=2\eta_q=0$, substituting
$[M_q]=d(1+s)+\eta_q$ gives $\eta_q\eta_{q'}=0$.
This proves all of (C.1.13).
\end{proof}

For $A=\mathbb Z_{(2)}$ or $\mathbb Z_2$, the index set in
(C.1.12) is countably infinite. The earlier map to $\mathbb Z/4$
detects those $q$ for which $q(t^2+t)$ is irreducible; the
$\epsilon$-Bockstein also detects the other irreducibles.
For example $q=u^3+u+1$ has trace-zero roots and
$q(t^2+t)$ splits into two cubics exchanged by $t\mapsto t+1$,
whereas $\chi_q^\epsilon(\eta_q)=1$. This explains why the
first Bockstein calculation alone could only prove a subgroup
injection and why the full calculation requires another invariant.

The finite certificate \texttt{code/C\_finish\_dualnumber\_ring.py}
constructs the integral Tate groups from actual kernels, images and
the modulo-four Bockstein, then computes the $\epsilon$-Bockstein
from the first-order matrix. Among $6736$ block extension candidates
on seven endpoint objects it checks all $197$ admissible extensions
for cyclic exactness, Bockstein compatibility and primary parity
additivity. It checks all $292$ two-cycle matrix pairs over
$\mathbb F_2$ in dimensions $(1,1),(1,2),(2,1),(2,2)$, lifting
all $2076$ invariant subspace pairs by integral even-entry gauges.
Finally it checks the seven nonzero irreducibles of degrees at most
four and all $49$ corresponding tensor identities (C.1.16).
The finite scope, matrices and runtime are recorded in
\texttt{logs/C\_finish\_dualnumber\_ring.json}; the general
classification rests on the proof above.
